\documentclass{article}

\usepackage{float}

\usepackage{enumitem}%многоуровневая нумерация в алгоритме
\usepackage[left=2.3cm,right=2.5cm,top=2.5cm,bottom=3.5cm]{geometry}
\usepackage{amsmath}
\usepackage{graphicx}
\usepackage[colorlinks=true, allcolors=blue]{hyperref}
\usepackage[T2A]{fontenc}
\usepackage[utf8]{inputenc}
\pagestyle{fancy} 
\usepackage[english,russian]{babel}
\begin{document}




\begin{titlepage}

  \begin{center}

    \vspace{6cm}
      Министерство образования Республики Беларусь\\[1em]
    Учреждение образования\\
    БЕЛОРУССКИЙ ГОСУДАРСТВЕННЫЙ УНИВЕРСИТЕТ \\
    ИНФОРМАТИКИ И РАДИОЭЛЕКТРОНИКИ\\[1em]

  \begin{minipage}{\textwidth}
      \begin{flushleft}
        \begin{tabular}{ l l }
          Факультет & Информационных технологий и управления\\
          Кафедра   & Интеллектуальных информационных технологий
        \end{tabular}
      \end{flushleft}
    \end{minipage}\\[1em]

    
    \vspace{5em}


    \textbf{РАСЧЕТНАЯ РАБОТА}\\
    {по дисциплине <<Представление и обработка информации в интеллектуальных системах>>}\\
    на тему\\
    \textbf{\large Решение теоретико-графовой задачи.Найти обхват графа}\\[1em]

    \vspace{10em}
    
    \begin{tabular}{ p{0.65\textwidth}p{0.25\textwidth} }
      Выполнил:& М.\,Д.~Дорош \\[1em]
      Студент группы& \\
      421703 & \\
      Проверил: & Н.\,В.~Зотов \\
     
    \end{tabular}
    
    \vfill
    {\normalsize Минск 2025}
  \end{center}
\end{titlepage}
\newpage 
\section{Формализация алгоритма расчетной работы (РР)}
\textbf{Задание РР} 

2.13 Найти обхват орграфа

\textbf{Список используемых понятий:}
\textbf{1.Граф }- математическая структура, состоящая из вершин и ребер, соединяющих некоторые пары вершин.
\begin{figure}[H]
    \centering
    \includegraphics[width=0.5\linewidth]{граф.jpg}
    \label{fig:enter-label}
\end{figure}

\textbf{2.Неориентированный граф} — это геометрическая фигура, 
состоящая из точек (вершин графа) и соединяющих их линий (рёбер графа).
\begin{figure}[H]
    \centering
    \includegraphics[width=0.5\linewidth]{неорграф.jpg}
    \label{fig:enter-label}
\end{figure}
\textbf{3.Ориентированный граф} — это геометрическая фигура, состоящая из точек (вершин графа) и соединяющих их линий (рёбер графа), причём эти линии имеют направление.

\begin{figure}[H]
    \centering
    \includegraphics[width=0.5\linewidth]{орграф.jpg}
    \label{fig:enter-label}
\end{figure}
\textbf{Алгоритм РР}

Инициализация переменной mingirth значением n + 1, где n — количество узлов в графе. Это больше, чем максимально возможная длина цикла в графе, поэтому mingirth будет обновляться при обнаружении более короткого цикла.
Для каждого узла i в графе (внешний цикл): a. Инициализация вектора visited размером n, заполненного значениями false. Этот вектор будет отслеживать, какие узлы были посещены во время обхода в ширину (BFS). b. Инициализация очереди q с начальным узлом i и расстоянием 0. c. Запуск BFS, пока очередь q не станет пустой: d. Извлечение из очереди текущего узла cur и расстояния dist. e. Для каждого столбца j в матрице matrix:
Если matrix[cur][j] == 1, это означает, что в этом столбце есть ребро.
Тогда код проверяет все узлы neighbor в этом столбце, где matrix[neighbor][j] == -1, что означает отсутствие ребра.
Если neighbor == i то это означает, что цикл найден, и mingirth обновляется до минимального значения между текущим mingirth и dist + 1 (длина цикла).
Если neighbor != i и visited[neighbor] == false, то neighbor добавляется в очередь q с расстоянием dist + 1.
После завершения BFS для всех узлов i проверяется значение mingirth:
Если mingirth осталось равным n + 1, это означает, что в графе не было найдено ни одного цикла, поэтому функция возвращает -1.
Иначе функция возвращает значение mingirth - минимальную длину цикла, найденную в процессе.

\textbf{Результат формализации алгоритма РР}

Пользователь задает ориентированный граф.
\begin{figure} [H]
    \centering
    \includegraphics[width=0.5\linewidth]{graph.png}
    \label{fig:enter-label}
\end{figure}


Выбирается вершина для начала обхода (допустим, A). Инициализируется переменная mingirth, которая будет хранить обхват орграфа, изначально равная n(количеству вершин графа)+1. Также создаётся вектор visited, который хранит посещённые вершины. И очередь q, которая хранит текущую вершину и расстояние.
\begin{figure}[H]
    \centering
    \includegraphics[width=0.5\linewidth]{wave1.png}
    \label{fig:enter-label}
\end{figure}
Используя поиск в ширину, мы переходим к вершине B, обновляя очередь q и сохраняя её в векторе visited.
\begin{figure}[H]
    \centering
    \includegraphics[width=0.5\linewidth]{wave2.png}
    \label{fig:enter-label}
\end{figure}
\textbf{4.Обхват гарфа} - длина наименьшего цикла в графе.
Продолжаем обход в ширину и все вышеперечисленные процессы.
\begin{figure}[H]
    \centering
    \includegraphics[width=0.5\linewidth]{wave3.png}
    \label{fig:enter-label}
\end{figure}
Дойдя до вершины D, мы используем алгоритм обхода в ширину и получаем, что смежная вершина A является начальной вершиной. И получаем конечную дистанцию.
\begin{figure}[H]
    \centering
    \includegraphics[width=0.5\linewidth]{wave4.png}
    \label{fig:enter-label}
\end{figure}

Прописываем алгоритм сравнения текущей дистанции с переменной mingirth для получения минимального значения. И затем присваиваем mingirth минимальное значение.
\begin{figure}[H]
    \centering
    \includegraphics[width=0.5\linewidth]{wave5.png}
    \label{fig:enter-label}
\end{figure}

Потворяя данный алгоритм для всех вершин получим обхват орграфа.
Пример входных и выходных конструкций:
Входная:
\begin{figure}[H]
    \centering
    \includegraphics[width=0.5\linewidth]{test1.1.png}
    \label{fig:enter-label}
\end{figure}


Выходная:
\begin{figure}[H]
    \centering
    \includegraphics[width=0.5\linewidth]{test1.2.png}

    \label{fig:enter-label}
\end{figure}

Входная:

\begin{figure}[H]
    \centering
    \includegraphics[width=0.5\linewidth]{test2.1.png}
    \label{fig:enter-label}
\end{figure}

Выходная:

\begin{figure}[H]
    \centering
    \includegraphics[width=0.5\linewidth]{test2.2.png}
    \label{fig:enter-label}
\end{figure}

Входная:

\begin{figure}[H]
    \centering
    \includegraphics[width=0.5\linewidth]{test3.1.png}

    \label{fig:enter-label}
\end{figure}

Выходная:
\begin{figure}[H]
    \centering
    \includegraphics[width=0.5\linewidth]{test3.2.png}
    \label{fig:enter-label}
\end{figure}
\textbf{Вывод}
Формализовал алгоритм получения обхвата графа. Улучшил свои навыки работы с KBE.

\newpage
\textbf{Список литературы}
[1] Оре О. Теория графов. – 2-е изд. – М.: Наука, 1980. – С. 336.

[2] Кормен Т. Х. и др. Часть VI. Алгоритмы для работы с графами // Алгоритмы: построение и анализ = Introduction to Algorithms. – 2-е изд. – М.: Вильямс, 2006. – С. 1296.

[3] Харари, Ф. Теория графов / Ф. Харари / Пер. с англ. и предисл. В.П. Козырева. Под ред. Г.П. Гаврилова. Изд. 2-е. – М.: Едиториал УРСС, 2003. – 269 с.

[4] Нечипуренко, М. И. Алгоритмы и программы решения задач на графах и сетях / М. И. Нечипуренко, В. К. Попков, С. М. Майнагашев и др. – Новосибирск: Наука. Сиб. отд-ние, 1990. – 515 с.
\end{document}